% 圆锥 P-底面圆O（2011 湖南理19）：C为弧AB中点，D为AC中点
\begin{tikzpicture}[x=1cm,y=1cm,line width=.6pt,font=\normalsize,>=Stealth]
  \coordinate (P) at (0,5);
  \coordinate (O) at (0,0);
  \coordinate (A) at (-3,0);
  \coordinate (B) at (3,0);
  \coordinate (C) at (1.5000000000,0.6062177826);
  \coordinate (D) at (-0.7500000000,0.3031088913);
  \draw[dashed] (B) arc (0:180:3 and 0.7);
  \draw (A) arc (180:360:3 and 0.7);
  \draw (P) -- (A);
  \draw (P) -- (B);
  \draw[dashed] (P) -- (O);
  \draw[dashed] (A) -- (B);
  \draw[dashed] (A) -- (C);
  \draw[dashed] (P) -- (D);
  \draw[dashed] (P) -- (C);
  \draw[dashed] (O) -- (D);
  \node[above,yshift=3.0pt] at (P) {$P$};
  \node[left,xshift=-3pt] at (A) {$A$};
  \node[right,xshift=3pt] at (B) {$B$};
  \node[below=2pt] at (O) {$O$};
  \node[above right,xshift=3pt,yshift=3pt] at (C) {$C$};
  \node[above left=2pt,xshift=-1pt,yshift=7pt] at (D) {$D$};
\end{tikzpicture}
